\chapter{Bleeding After Vaginal Sex (Postcoital Bleeding)}

\section*{Definition}
Vaginal bleeding or spotting occurring after sexual intercourse, ranging from light spotting to heavier flow, originating from the cervix, vagina, or uterus.

\section*{When to See a Doctor}

\begin{itemize}
 \item Arrange assessment for any unexplained bleeding after sex, especially if it recurs, is heavy, occurs after menopause, or is associated with pelvic pain, discharge, or pain during sex.
 \item Seek urgent care for heavy bleeding, faintness, severe pelvic pain, or bleeding after a missed period, because ectopic pregnancy must be excluded.
\end{itemize}

\section*{How It Develops}
Postcoital bleeding can arise from the cervix, vagina, or uterus. Common causes include cervical ectropion, cervicitis, cervical polyps, vaginal dryness, and small tears from friction or trauma. Less commonly, it is associated with cervical or endometrial precancer/cancer, polyps or fibroids, pregnancy-related conditions, or medicines that increase bleeding.

\section*{Causes}
\begin{itemize}
 \item Cervicitis (chlamydia, gonorrhea, trichomonas, HSV)
 \item Cervical ectropion (ectopy)
 \item Cervical or endometrial polyps
 \item Cervical precancer or cancer (uncommon, but important to exclude)
 \item Vaginal atrophy (genitourinary syndrome of menopause)
 \item Insufficient lubrication
 \item Vaginal or cervical trauma/laceration
 \item Endometritis or pelvic inflammatory disease
 \item Pregnancy (cervical friability)
 \item Coagulation disorders or anticoagulant therapy
 \item Hormonal contraception or other hormonal changes
\end{itemize}

\section*{Related Symptoms}
\begin{itemize}
 \item Vaginal bleeding between periods
 \item Pelvic pain
 \item Vaginal discharge
 \item Painful sex
 \item Abnormal Pap smear
\end{itemize}


\section*{Medical Evaluation}
\begin{itemize}
 \item Your doctor will perform a speculum examination to visualise the cervix and vagina, identifying ectropion, polyps, lacerations, or lesions.
 \item A pregnancy test is performed when pregnancy is possible. Tests for sexually transmitted infections are selected according to symptoms, sexual history, and local guidance.
 \item Cervical screening is kept up to date, but a screening test is not a diagnostic test and should not replace examination of symptomatic bleeding.
 \item Depending on findings, further assessment may include ultrasound, colposcopy, or hysteroscopy; referral is arranged for an abnormal cervix, abnormal screening result, persistent symptoms, or suspected cancer.
\end{itemize}

\section*{Treatment Options}
\begin{itemize}
 \item Cervical ectropion is managed conservatively or with cautery if bleeding is persistent and bothersome.
 \item Antibiotic or antiviral therapy is prescribed for a diagnosed infection according to current local guidelines and resistance patterns; partners may also need testing or treatment.
 \item Vaginal moisturisers, lubricants, or locally prescribed vaginal oestrogen may help genitourinary menopause symptoms after appropriate assessment and discussion of risks and benefits.
 \item Polypectomy (cervical or endometrial) is performed for polyp-related bleeding, and loop electrosurgical excision (LEEP) is used for high-grade cervical dysplasia.
\end{itemize}

\section*{Self-Care and Home Management}
\begin{itemize}
 \item If sex causes pain or bleeding, pause penetration until assessed and symptoms settle; seek advice rather than assuming the cause is friction.
 \item Use a panty liner to monitor the amount and colour of bleeding — note whether it is bright red, dark, or brown-tinged.
 \item Apply a water-based lubricant during intercourse to reduce friction if vaginal dryness or atrophy is contributing to bleeding.
 \item Avoid douching or using vaginal products (creams, sprays, wipes) which can irritate the cervix and vaginal mucosa.
\end{itemize}

\section*{References}
\begin{thebibliography}{99}
\bibitem{bleeding_after_vaginal_sex:ref1} Tarney CM, Han J. ``Postcoital bleeding: a review of the literature.'' \textit{J Womens Health}. 2014;23(4):307-312.
\bibitem{bleeding_after_vaginal_sex:ref2} Luesley DM, Sidhu HK. ``Postcoital bleeding.'' \textit{BMJ}. 2009;338:b692.
\bibitem{bleeding_after_vaginal_sex:ref3} Saccardi C, Spagnol G, et al. ``Postcoital bleeding: an update on evaluation and management.'' \textit{Minerva Obstet Gynecol}. 2021;73(4):448-456.
\bibitem{bleeding_after_vaginal_sex:ref4} Berek JS, Berek DL, et al. \textit{Berek \& Novak\'s Gynecology}, 16th ed. Wolters Kluwer, 2019.
\bibitem{bleeding_after_vaginal_sex:ref5} UpToDate. ``Causes and evaluation of postcoital bleeding.'' Wolters Kluwer, 2023.
\end{thebibliography}
